% Part V: carriers, the room growing, and the assembly.
% ASCII only. Layers 6, 7, 5 and 1 of the grid, then the assembly of the chain.
\newpage
\proofpart{5}
\pagetitle{25 / PART V. CARRIERS, THE ROOM GROWING, AND THE ASSEMBLY}{From the small room to the continuum}

This part covers the carrier and the switch (layers 6, 7), topology (layer 5) and the room's size (layer 1), then
assembles the chain. Most of it is conjecture, each stated so that it can fail.

\textbf{Notation.} $H_{\rm src}(s)$ is the sourced reduced Hamiltonian of page 11 on $\mathbb R^{3\times3}$, at fixed
$t$ and $g$, with sources $\mathrm{src}=(a,b,b)$, $a=(1-s)(t-2s)$, $b=(1-s)((1-s)/t-s)$, on the window
$1<s<(t+2)/(t+1)$; the \emph{charge} is the Berger $U(1)$, rotating frame indices 2 and 3. $\Delta_H(s)$ is the
excitation energy of the lowest charge-$\pm2$ level (Helike's channel) and $\Delta_C(s)$ that of the second charge-0
level (the competitor's), both above the ground energy, in the colour-singlet sector. A \emph{carrier} is followed
across $s$ by its physical overlap, degenerate partners together; its \emph{own-block share} $\sigma$ is the fraction
of its vacuum-weighted Dirichlet energy spent in its own normal-mode coordinates (M8). Theorems 8 and 9 are those of
my bundle-framing note (Chern--Simons number on the line; vacua and background of the reduced model).

\medskip\eyebrow{LAYERS 6 AND 7: THE CARRIER AND THE SWITCH}

\begin{proposition}[the switch is an exact crossing]\label{prop:V-crossing}
\status{PROVED HERE}
Fix $t>1$, $g>0$. (i) $H_{\rm src}(s)$ commutes with the charge rotation and with colour $SO(3)$ for every $s$.
(ii) On the window, $\{H_{\rm src}(s)\}$ is a self-adjoint holomorphic family of type (B) in $s$; in each charge
sector the eigenvalues, counted with multiplicity, are continuous in $s$ and are given locally by real-analytic
branches. (iii) The charge-$\pm2$ and charge-0 sectors are invariant, so $H_{\rm src}(s)$ has no matrix element
between Helike's channel and the competitor's: where $\Delta_C-\Delta_H$ changes sign, the levels cross exactly; there
is no avoided crossing. If $\Delta_C-\Delta_H$ has opposite signs at $s_1<s_2$, it vanishes somewhere in between.
\end{proposition}
\begin{proof}
(i) The Laplacian is invariant under orthogonal changes of rows and of columns. For $R,O\in SO(3)$, $\operatorname{cof}(Rq)=R\operatorname{cof}(q)$ and $\operatorname{cof}(qO)=\operatorname{cof}(q)O$.
Let $R$ be a rotation in the frame plane (2,3). Because $\omega_2=\omega_3$ and $\mathrm{src}_2=\mathrm{src}_3$,
$\omega\circ(Rq)=R(\omega\circ q)$ and $\sum_k\mathrm{src}_k\|(Rq)_k\|^2=\sum_k\mathrm{src}_k\|q_k\|^2$. So
$V(Rq)=V(q)$, and likewise $V(qO)=V(q)$. (ii) Only the sources depend on $s$, polynomially. Fix $s_0$. Theorem A gives $V_{s_0}\geq\frac{\kappa}{2}\|q\|^2-S^4/(2g^2)$ with $\kappa=\kappa(s_0)>0$, so, as quadratic forms,
$\|q\|^2\leq(2/\kappa)\bigl(H_{\rm src}(s_0)+S^4/(2g^2)\bigr)$. The perturbation
$H_{\rm src}(s)-H_{\rm src}(s_0)=-\frac12\sum_k(\mathrm{src}_k(s)-\mathrm{src}_k(s_0))\|q_k\|^2$ is therefore
form-bounded with relative bound at most $\max_k|\mathrm{src}_k(s)-\mathrm{src}_k(s_0)|/\kappa$, which tends to zero as
$s\to s_0$, also for complex $s$. So the family is holomorphic of type (B) (Kato, \emph{Perturbation Theory for Linear
Operators}, Ch.~VII), self-adjoint for real $s$; each sector has compact resolvent (Theorem A), so Rellich's theorem
gives analytic branches, and min--max gives continuity. (iii) follows from (i); the last claim is the intermediate
value theorem.
\end{proof}

So an ``identifiable crossing'' is two branches in different sectors meeting. Within a sector, tracking by overlap is
still needed: the competitor's branch can meet another charge-0 branch.

\begin{conjecture}[the switch and carrier persistence, M8]\label{conj:V-switch}
\status{CONJECTURE}
At $t=3$ and $g^2\in\{0.03,0.05,0.07\}$ there is exactly one $s^*$ in the registered $s$-range with
$\Delta_C(s^*)=\Delta_H(s^*)$, with $\Delta_C>\Delta_H$ for $s>s^*$ (Helike carries the lowest note) and
$\Delta_C<\Delta_H$ for $s<s^*$ (the competitor carries it), and $|s^*-s^*_{\rm pred}|\leq0.015$ with the registered
predictions $s^*_{\rm pred}=1.1674$, $1.1809$, $1.1931$. On each side of $s^*$ each carrier keeps own-block share
$\sigma\geq\frac12$ while it is its channel's lowest level, so that the M8 verdict is PERSISTS or HANDS OVER, never
DISSOLVES.
\end{conjecture}
\begin{sketch}
\emph{Step 1, the mechanism.} At weak coupling each gap is $\omega_X+c_Xg^2+O(g^4)$. At $t=3$, $s=1.15$ the harmonic
frequencies are $\omega_H=1.090107$ (doublet) and $\omega_C=1.106657$, and the exact $O(g^2)$ coefficients are
$c_H=-0.530$ and $c_C=-2.738$. At $g^2=0.002$, $0.01655-2.208\,g^2=0.0121$ matches the converged gaps $1.10114-1.08904$; at
$g^2=0.03$ it is negative at $s=1.15$, so the crossing moves to larger $s$. Solving
$\omega_C(s)-\omega_H(s)+(c_C(s)-c_H(s))g^2=0$ gives the registered $s^*_{\rm pred}$, and the tolerance $0.015$
absorbs $O(g^4)$, which largely cancels between the two levels.
\emph{Step 2, the anchor.} At the controls the ordering $\Delta_H<\Delta_C$ is proved at three of four points,
conditional on $S_C$ (Part III). That fixes the sign on the right of the crossing at weak coupling.
\emph{Step 3, measuring the crossing.} The registered grid (gate57) is UNRESOLVED: the competitor had neither
converged at $M_q=10$ nor a clean bare identity. The dressed chart (\S5.10) repairs the identity
($\tilde w=0.957$ against $0.914$ bare at the smoke point $s=1.18$, $g^2=0.03$) but not the rate. The unfolded chart
$P=\sqrt G$ (\S5.11) is registered to repair the rate; an exploratory prototype gave a competitor step of $0.0004$ from
$M_q$ 6 to 8 at the smoke point, seen before registration, so it decides nothing, and the grid is held. With converged
gaps at the grid points,
Proposition~\ref{prop:V-crossing} turns a sign change between adjacent points into an exact crossing.
\emph{Step 4, identity.} At the controls both carriers are ALIGNED, with smallest share $0.9956$. The share is
continuous in $s$ while the carrier's branch is isolated in its sector, so it can fall below $\frac12$ only through a
within-sector near-crossing, which M8 detects by refinement.
\emph{Step 5, the fast fibre.} The full theory dresses each slow mode: to first order
$\delta\omega_k=g^2(u_k\cdot\operatorname{Hess}_FE_0^{\rm fast}\cdot u_k)/\omega_k$, up to the scheme freedom
$E_0\to E_0+\kappa\,dC$. At the background the ordering survives if the difference of the two shifts is below the margin
there ($0.01655$ harmonic, $0.0121$ converged at $g^2=0.002$); elsewhere the shift moves $s^*$ to first order (Part IV). At the background this curvature is \MthreeBackground. The conjecture is stated for the reduced model and
does not depend on that value; its extension to the full theory does.
\emph{Where it could fail.} The competitor's $O(g^2)$ coefficient is five times Helike's, so its higher orders may not
cancel as they did at small $g$; its level was not converged on the grid, and no grid value is released. A converged
crossing outside tolerance, or a share below $\frac12$ without a verified crossing, refutes it.
\end{sketch}

\begin{conjecture}[the glueball link]\label{conj:V-glueball}
\status{CONJECTURE}
$s^*(g^2)$ is increasing in $g^2$, and stays increasing along the M2 sweep $g^2=g^2(R)$. The region $s<s^*$ where the
charge-0 carrier is lowest therefore grows with the room. In infinite volume it becomes the ordering
$m(0^{++})<m(2^{++})$ of the lattice glueball spectrum (Morningstar and Peardon, Phys.\ Rev.\ D 60 (1999) 034509).
\end{conjecture}
\begin{sketch}
At first order $s^*$ solves $\delta\omega(s)=-(c_C-c_H)g^2$, $\delta\omega=\omega_C-\omega_H$, so
$ds^*/dg^2=-(c_C-c_H)/\delta\omega'(s^*)>0$ when $c_C<c_H$ (as at the controls) and $\delta\omega'>0$. A larger $g^2$ is a
larger room: $R\Lambda=\exp(-8\pi^2/(b_0g^2))$ increases with $g^2$. Charge 0 contains the scalar; the charge-$\pm2$ doublet is the
$m=\pm2$ part of the spin-2 part of $G=qq^T$, a quintet on the round sphere. \emph{Where it could fail.} First order
is not a proof at $g^2=0.07$, and no converged grid value yet tests the rise. The step from $g^2$ to $R$ needs the source bridge
(Conjecture~\ref{conj:V-bridge}). The labels must be continued to $J^{PC}$ past the volume where the reduction
fails, and the van den Heuvel crossing may be a truncation artefact (note \S7). I quote no mass ratios.
\end{sketch}

\medskip\eyebrow{LAYER 5: TOPOLOGY ON THE MAURER--CARTAN LINE}

\begin{proposition}[the barrier and the tunnelling length on the line]\label{prop:V-agmon}
\status{PROVED HERE}
With the sources off, on the line $q=\lambda q_{\rm MC}$, $q_{\rm MC}=\operatorname{diag}(1/\sqrt t,1,1)/g$:
(i) the barrier is the sphaleron value $V(\tfrac12q_{\rm MC})=(1+2/t)/(32g^2)$; (ii) the Agmon length of the segment
$0\leq\lambda\leq1$, for $-\frac12\Delta_q+V$, is
\[
 L(t)=\int_0^1\sqrt{2V(\lambda q_{\rm MC})}\,\|q_{\rm MC}\|\,d\lambda=\frac{\sqrt{(t+2)(2t+1)}}{6\,t\,g^2},
\]
an upper bound for the Agmon distance between the two vacua; (iii) at $t=1$ the Agmon distance in $\mathbb R^9$ from
$0$ to the orbit $SO(3)/g$ equals $1/(2g^2)=L(1)$ exactly.
\end{proposition}
\begin{proof}
(i) and (ii): Theorem 9 gives $V(\lambda q_{\rm MC})=(1+2/t)(\lambda-\lambda^2)^2/(2g^2)$, maximal at
$\lambda=\frac12$. Also $\|q_{\rm MC}\|^2=(1/t+2)/g^2$. Then
$\sqrt{2V}\,\|q_{\rm MC}\|=\sqrt{(t+2)(2t+1)}\,(\lambda-\lambda^2)/(tg^2)$ and $\int_0^1(\lambda-\lambda^2)d\lambda=1/6$.
Any path bounds the infimum. (iii) At $t=1$, $V=\frac12\|\nabla W\|^2$ with $W=\frac12T_q-g\det q$ (page 10). For any
path $\gamma$ from $0$ to $O/g$, $\int\sqrt{2V}\,|d\gamma|=\int\|\nabla W\|\,|d\gamma|\geq\int\nabla W\cdot d\gamma
=W(O/g)-W(0)=3/(2g^2)-1/g^2=1/(2g^2)$, and $L(1)=\sqrt9/(6g^2)$ attains it.
\end{proof}

\begin{proposition}[$\theta$ is a Bloch phase]\label{prop:V-theta}
\status{CONDITIONAL ON a self-adjoint Hamiltonian on gauge-orbit space that commutes with large gauge transformations}
Let $\mathcal T$ implement a degree-1 gauge transformation on states invariant under degree-0 ones. Then $H$
decomposes over the spectrum $e^{i\theta}$ of $\mathcal T$, and $\theta$ labels superselection sectors (Jackiw and Rebbi, Phys.\ Rev.\ Lett.\ 37
(1976) 172; Callan, Dashen and Gross, Phys.\ Lett.\ B 63 (1976)).
\end{proposition}
\begin{proof}
$\mathcal T$ is unitary and commutes with $H$; apply the spectral theorem to the commuting pair.
\end{proof}

\begin{remark}[what the line is, and a correction]
\status{PROVED HERE}
$\mathcal T$ shifts $CS$ by $1$ and carries the winding-0 vacuum to the winding-1 vacuum, the two ends of the segment
$0\leq\lambda\leq1$ (Theorem 8). The periodic chain of vacua is the orbit of this segment under $\mathcal T^n$, not the
continuation of $\lambda$ past $1$: by Theorems 8 and 9, $CS(2)=-4\neq CS(1)+1$ and
$V(2q_{\rm MC})=2(1+2/t)/g^2\neq0$. ``$\theta$ is the Bloch
phase of $\lambda\mapsto\lambda+1$'' is shorthand for this. The equivalence of the two ends needs the fast fibre: in the
reduced model $0$ is a point and the other vacuum an $SO(3)$ orbit, but the full slice spectrum at $\lambda=1$ equals
the one at $\lambda=0$ with multiplicities, its Peter--Weyl sectors reassembled (Part IV,
Proposition~\ref{p4prop:flat}); the regulated sums agree to $10^{-15}$ ($t=1,3$; computed).
\end{remark}

\begin{conjecture}[the $\theta$-dependence of the gap is exponentially small]\label{conj:V-theta}
\status{CONJECTURE}
On the Berger sphere at weak coupling, $|\Delta(\theta)-\Delta(0)|\leq C\,e^{-S_{\rm inst}}$, with
$S_{\rm inst}=8\pi^2/g^2$ in the normalization $\frac1{4g^2}\int F^a_{\mu\nu}F^{a\mu\nu}$ and $C$ at most polynomial
in $1/g^2$.
\end{conjecture}
\begin{sketch}
In tight binding each level of a vacuum spreads into a band $E(\theta)=E_0-2\tau\cos\theta$,
$\tau\sim e^{-d}$ with $d$ the Agmon distance through the sphaleron; the gap, a difference of band energies, moves by
$O(\tau)$. Proposition~\ref{prop:V-agmon}
gives $d$ exactly at $t=1$ in model units, $1/(2g^2)$, and an upper bound for $t>1$. Not checked: that the map from the model's $g$ to
the Yang--Mills coupling (field map $F=2gq$, volume) carries $1/(2g^2)$ to $8\pi^2/g^2$, and that the minimizing path
for $t>1$ is the line. \emph{Where it could fail:} the fast-fibre determinant in the prefactor; excited bands wider
than the ground band. Nothing in the chain uses this at $\theta=0$.
\end{sketch}

\medskip\eyebrow{LAYER 1: THE ROOM GROWS}

\begin{conjecture}[M2, the infrared half]\label{conj:V-M2}
\status{CONJECTURE}
Let $\Delta(R)$ be the gap of $SU(2)$ Yang--Mills on $\mathbb R\times S^3_R$ (fixed squashing $t$), and let $g^2(R)$
be given by $R\Lambda=\exp(-8\pi^2/(b_0g^2))$, $b_0=22/3$ (Part 0's convention; the note writes
$\exp(-1/(2b_0g^2))$ with its $b_0=11N/(48\pi^2)$). Then $\inf_{R>0}\Delta(R)/\Lambda>0$, and the
connected correlations of local gauge-invariant observables decay in Euclidean time at least like
$e^{-\Delta_{\min}|\tau|}$ with $\Delta_{\min}>0$ uniform in $R$.
\end{conjecture}
\begin{sketch}
\emph{Small room.} As $R\to0$, $g\to0$ and, granted M3 (Part IV), $\Delta(R)R$ tends to the lowest frequency of the
quadratic theory on the unit sphere, a positive constant, so
$\Delta/\Lambda=(\Delta R)/(R\Lambda)\to\infty$. The scale is consistent: the log-divergent coefficient of the fast
zero-point sum equals $(b_0/8\pi^2)\operatorname{Vol}(S^3_t)$ times the classical energy density, to $10^{-4}$ at $\lambda\in\{\frac18,\dots,\frac78\}$ (computed, Part IV). \emph{Middle.} If the theory exists at each $R$ with a positive gap
depending continuously on $R$, then $\Delta/\Lambda$ is positive on every compact interval of $R$. So everything
reduces to $\liminf_{R\to\infty}\Delta(R)/\Lambda>0$. \emph{Large room.} At strong bare coupling the lattice theory
has a gap and clustering by a convergent expansion (Osterwalder and Seiler, Ann.\ Phys.\ 110 (1978) 440; Part I). The
continuity argument along scale is the inherited note (note \S4.9): follow the lowest level from the strong-coupling
Casimir end to the weak-coupling small-room end with a traceable carrier, and show it never reaches zero in units of
$\Lambda$. \emph{The key estimate} would be a lower bound on the gap of the block-spin effective theory at the scale
where its coupling becomes of order one, uniform in the volume: that is, confinement. \emph{Where it fails.} Nothing
I have controls the strongly coupled scales, and a gap that is positive at every finite $R$ need not survive in units
of $\Lambda$ (the analogue for cutoffs: three-dimensional compact $U(1)$ is gapped at every cutoff, yet $m/\sqrt\sigma\to0$;
note \S4.9). This is the wall.
\end{sketch}

\begin{conjecture}[the source bridge: an open condition]\label{conj:V-bridge}
\status{CONJECTURE}
The spectrum of $H_{\rm src}(s)$ determines, by a controlled limit or a matching, the low spectrum of the unsourced
reduction, and through M3 that of the full theory on $S^3_R$.
\end{conjecture}
\begin{sketch}
The sources vanish exactly at the window's edge: $a=b=0$ at $s=1$, where the background is the winding-1 vacuum
$q_{\rm MC}$ (Theorem 9). The route: follow every reading as $s\to1^+$ and control the limit. \emph{Where it could
fail:} the unsourced vacuum spreads over both wells, so the single-well picture may not survive, and a switch at
$s^*>1$ may have no counterpart at $s=1$. A running coupling assigned to a radius does not supply this. Until it is
crossed, M2 is a statement about the sourced reduction.
\end{sketch}

\medskip\eyebrow{THE ULTRAVIOLET: EXISTENCE}

\begin{theorem}[partial ultraviolet results]\label{thm:V-uv}
\status{PROVED, CITED: Osterwalder and Seiler 1978; Ba{\l}aban, 1980s; Magnen, Rivasseau and S\'en\'eor 1993}
(i) The Wilson lattice measure is reflection positive (Osterwalder and Seiler, Ann.\ Phys.\ 110 (1978) 440). (ii)
Ba{\l}aban's renormalization-group analysis bounds the effective actions of four-dimensional lattice Yang--Mills in a
finite volume uniformly in the lattice spacing (Ba{\l}aban, Commun.\ Math.\ Phys., series of papers in the 1980s).
(iii) The Schwinger functions of pure Yang--Mills in four dimensions exist with the ultraviolet cutoff removed and a
fixed infrared cutoff (Magnen, Rivasseau and S\'en\'eor, Commun.\ Math.\ Phys.\ 1993).
\end{theorem}

\begin{proposition}[reflection positivity survives limits]\label{prop:V-rp}
\status{PROVED HERE}
Let $\mu_n$ be Wilson lattice measures and $\Theta$ a reflection in a hyperplane that is a lattice hyperplane for
every $n$. If $\langle F\rangle_{\mu_n}\to\langle F\rangle$ for every bounded local observable $F$, then
$\langle\Theta(F)\,F\rangle\geq0$ for every bounded local $F$ supported on the positive side.
\end{proposition}
\begin{proof}
$\Theta(F)F$ is bounded and local, so $\langle\Theta(F)F\rangle=\lim_n\langle\Theta(F)F\rangle_{\mu_n}$, and each term is
$\geq0$ by Theorem~\ref{thm:V-uv}(i). A limit of non-negative numbers is non-negative.
\end{proof}

\begin{conjecture}[ultraviolet existence with the Osterwalder--Schrader axioms]\label{conj:V-uv}
\status{CONJECTURE}
For $SU(2)$ in a finite volume, the lattice Schwinger functions of gauge-invariant observables converge as $a\to0$
with $g^2(a)$ running as asymptotic freedom dictates, and the limit satisfies the Osterwalder--Schrader axioms with
the regularity bounds.
\end{conjecture}
\begin{sketch}
(1) \emph{Reflection positivity} passes to the limit by Proposition~\ref{prop:V-rp} (for reflections that stay
lattice hyperplanes; the others need (3)). (2) \emph{Limits}: uniform
bounds of Ba{\l}aban type give tightness, hence convergent subsequences; uniqueness needs more. (3) \emph{Euclidean
invariance}: gauge-invariant hypercubic operators that break $O(4)$ first appear at dimension 6, so they are irrelevant
and vanish like $a^2$ up to logarithms (Symanzik's argument); rigour needs the multiscale control of (2).
\emph{Where it could fail}: in gauge-invariant multiscale control itself (no global gauge choice, Singer).
\end{sketch}

\medskip\eyebrow{THE ASSEMBLY}

\begin{proof}[Proof sketch of the Main Conjecture from Parts I--V \textmd{(my best guess; not a proof)}]
Each step names what it uses and that result's status.
\begin{enumerate}[leftmargin=18pt,itemsep=2pt,topsep=2pt]
\item \emph{The lattice object.} A gauge-invariant, reflection-positive measure on lattice connections modulo
gauge, read by Wilson loops (Part I). \textbf{Proved; proved, cited.}
\item \emph{The note at strong coupling.} Gap and clustering at strong bare coupling (Osterwalder--Seiler); the
leading note is the Casimir $C_2(\frac12)=\frac34$ on the smallest loop; in two dimensions the area law with that
Casimir is exact (Part II). \textbf{Proved, cited; proved for the electric Hamiltonian (Part I) and in 2-D (Part II);
conjecture as the leading term of the full lattice gap.}
\item \emph{The small room, reduced.} Discrete spectrum of the sourced reduction (Theorem A); Helike is the classical
floor on Berger spheres (\emph{By Enough}); her quantum level lies below the competitor's at three of four controls
(Part III). \textbf{Proved; proved, cited; conditional on $S_C$.}
\item \emph{The small room, full.} The fast fibre only dresses the slow modes, uniformly in the cutoff (M3, Part IV).
Its divergence is the one-loop running; the two vacua are spectrally equivalent. \textbf{Conjecture; the divergence
computed, the equivalence proved (Part IV).}
\item \emph{Carrier and switch.} Proposition~\ref{prop:V-crossing} \textbf{(proved here)};
Conjectures~\ref{conj:V-switch}, \ref{conj:V-glueball} \textbf{(conjecture)}. \emph{Topology.}
Propositions~\ref{prop:V-agmon} \textbf{(proved here)}, \ref{prop:V-theta} \textbf{(conditional)};
Conjecture~\ref{conj:V-theta} \textbf{(conjecture; not used at $\theta=0$)}.
\item \emph{The source bridge.} Conjecture~\ref{conj:V-bridge}. \textbf{Open.}
\item \emph{The room grows.} Conjecture~\ref{conj:V-M2} (M2). \textbf{Conjecture: the wall.}
\item \emph{Existence.} Theorem~\ref{thm:V-uv} \textbf{(proved, cited; partial)} and Conjecture~\ref{conj:V-uv}
\textbf{(conjecture)}, in each finite volume.
\item \emph{The limit.} Grant 4 and 8 at each $R$, 6 to connect them to the reduction, and the uniform clustering of 7.
A bound $|\langle A\,B(\tau)\rangle_{\rm c}|\leq C_{A,B}\,e^{-\Delta_{\min}|\tau|}$ that is uniform in $R$ passes to
any limit $R\to\infty$, and reflection positivity passes with it (Proposition~\ref{prop:V-rp}). Osterwalder--Schrader reconstruction then gives a
Hilbert space and $H\geq0$ with no spectrum in $(0,\Delta_{\min})$ on the subspace generated by local observables:
the gap, $\Delta=c\Lambda$. \emph{Euclidean invariance, uniqueness of the infinite-volume limit and
non-triviality are further conditions, not supplied by any step above.}
\item \emph{Inheritance.} If, in addition, step 5's carrier conjectures hold along the M2 sweep and the carrier is traced from step 2 to step 3
(M9, on the lattice), the gap is the inherited note with an identified carrier. \textbf{Conjecture.}
\end{enumerate}
\emph{Which single failures break it.} This route to the gap (step 9) breaks if any of 4, 6, 7 or 8 fails. It does not break if
$S_C$ fails, if the switch is misplaced, or if the carrier dissolves: those break only the identified carrier and the
inheritance (step 10). Topology breaks nothing at $\theta=0$. Step 7 is the one I expect to be hardest, since it is the
infrared half of the Clay problem itself.

\emph{What the GIGI construction can and cannot test.} It can store each reading as a section and recompute it from
point queries; check planted exact answers (the 2-D area law, gauge invariance to $10^{-15}$, the centre vortex, the
closed forms on the Maurer--Cartan line, the Hessian slice spectrum); hold a carrier's identity across $s$, cutoffs
and charts, and refuse a match that is only a similar number. It cannot prove a bound, take $a\to0$ or $R\to\infty$,
or say anything in four dimensions beyond the finite computations it stores. An agreement it records is evidence, not
a proof.
\end{proof}

This is a complete construction, unvalidated, and overall unproven.
